\documentclass[aspectratio=169,10pt]{beamer}
% Shared preamble for the Algorithms and Programming Concepts slides (Beamer, Metropolis)
\usetheme[progressbar=frametitle,numbering=fraction,background=light,block=fill,sectionpage=none]{metropolis}
\usepackage[T1]{fontenc}
\usepackage{FiraSans}
\usepackage[scaled=.85]{FiraMono}
\usepackage{amsmath,amssymb}
\usefonttheme[onlymath]{serif}
\usepackage{booktabs,array,tabularx}
\usepackage{tikz,pgfplots}
\usepackage{xcolor}
\usepackage{listings}
\usepackage{adjustbox}
\usepackage{fontawesome5}
\hyphenpenalty=10000\exhyphenpenalty=10000
\newenvironment{sfig}[1][0.80]{\begin{adjustbox}{max width=\linewidth,max totalheight=#1\textheight,center}}{\end{adjustbox}}
\usetikzlibrary{arrows.meta,positioning,shapes.geometric,shapes.misc,fit,calc,backgrounds,decorations.pathreplacing,patterns,matrix,chains}
\pgfplotsset{compat=1.17}

\definecolor{tdInk}{HTML}{1D2330}
\definecolor{tdMuted}{HTML}{5D6678}
\definecolor{tdLine}{HTML}{D9D5CA}
\definecolor{tdPaper}{HTML}{FBFAF7}
\definecolor{tdBlue}{HTML}{2457D6}
\definecolor{tdGreen}{HTML}{17855A}
\definecolor{tdAmber}{HTML}{B7791F}
\definecolor{tdRed}{HTML}{C8362F}
\definecolor{tdPurple}{HTML}{7A4FD1}
\setbeamercolor{normal text}{fg=tdInk,bg=tdPaper}
\setbeamercolor{alerted text}{fg=tdRed}
\setbeamercolor{example text}{fg=tdGreen}
\setbeamercolor{frametitle}{bg=tdInk,fg=white}
\setbeamercolor{progress bar}{fg=tdBlue,bg=tdLine}
\setbeamercolor{title separator}{fg=tdBlue}
\setbeamercolor{block title}{bg=tdBlue!12,fg=tdInk}
\setbeamercolor{block body}{bg=tdBlue!5,fg=tdInk}
\setbeamercolor{block title alerted}{bg=tdRed!12,fg=tdRed!80!black}
\setbeamercolor{block body alerted}{bg=tdRed!5,fg=tdInk}
\setbeamercolor{block title example}{bg=tdGreen!14,fg=tdGreen!70!black}
\setbeamercolor{block body example}{bg=tdGreen!6,fg=tdInk}
\setbeamerfont{frametitle}{size=\large,series=\bfseries}
\makeatletter\setlength{\metropolis@frametitle@padding}{1.15ex}\makeatother
\setbeamertemplate{navigation symbols}{}
\setbeamerfont{block title}{size=\small}

% ---- code listings (Python 3.14) ----
\lstdefinestyle{py}{language=Python,
  basicstyle=\ttfamily\fontsize{8.4}{10.1}\selectfont\color{tdInk},
  keywordstyle=\color{tdBlue}\bfseries, commentstyle=\color{tdMuted}\itshape, stringstyle=\color{tdGreen},
  morekeywords={match,case,type,self,True,False,None},
  emph={print,len,range,sorted,min,max,sum,enumerate,zip,heappush,heappop,cache},emphstyle=\color{tdPurple},
  showstringspaces=false, columns=fullflexible, keepspaces=true, frame=single, rulecolor=\color{tdLine},
  backgroundcolor=\color{white}, xleftmargin=2pt, xrightmargin=2pt, aboveskip=2pt, belowskip=0pt, framesep=3pt}
\lstdefinestyle{out}{basicstyle=\ttfamily\fontsize{8}{9.6}\selectfont\color{tdInk}, columns=fullflexible, keepspaces=true,
  frame=single, rulecolor=\color{tdGreen!60}, backgroundcolor=\color{tdGreen!6}, xleftmargin=2pt, xrightmargin=2pt, aboveskip=2pt, belowskip=0pt, framesep=3pt}
% \pycode[size]{snippet} shows snippets/<snippet>.py ; \pyout[size]{snippet} shows snippets/<snippet>.out
\newcommand{\pycode}[2][\relax]{{#1\lstinputlisting[style=py]{snippets/#2.py}}}
\newcommand{\pyout}[2][\relax]{{#1\lstinputlisting[style=out]{snippets/#2.out}}}

% ---- tikz styles ----
\tikzset{>={Stealth[length=2.2mm]},
  card/.style={rounded corners=4pt,draw=#1!70!black,fill=#1!10,align=center,inner sep=5pt,font=\small}, card/.default=tdBlue,
  pill/.style={rounded corners=7pt,draw=none,fill=#1,text=white,inner sep=3pt,font=\scriptsize\bfseries,align=center}, pill/.default=tdBlue,
  arr/.style={->,thick,draw=tdMuted},
  cell/.style={draw=tdInk!70,minimum width=7.5mm,minimum height=7.5mm,inner sep=0pt,font=\small\ttfamily,fill=white},
  tn/.style={circle,draw=tdInk!80,fill=white,minimum size=7mm,inner sep=0pt,font=\small\bfseries}}
\newcommand{\key}[1]{\textbf{\color{tdBlue}#1}}
\newcommand{\bad}[1]{\textbf{\color{tdRed}#1}}
\newcommand{\good}[1]{\textbf{\color{tdGreen!70!black}#1}}
\newcommand{\apcfooter}[1]{\setbeamertemplate{footline}{\hfill\tiny\color{tdMuted}Algorithms and Programming Concepts \,|\, #1 \,|\, BUV\hspace{1.6em}\insertframenumber/\inserttotalframenumber\hspace{0.8em}\vspace{4pt}}}
\newcommand{\apctitle}[2]{\title{#1}\subtitle{#2}\author{Dr.\ Quang-Vinh Dang\\ {\small Senior Lecturer (Assoc.\ Prof)}}\institute{British University Vietnam}\date{Algorithms and Programming Concepts \,|\, Year 1 \,|\, Python 3.14}}

\apcfooter{Big-O and complexity}
\apctitle{Big-O and complexity}{How the cost of an algorithm grows with the size of its input}

\begin{document}
\begin{frame}[plain]\titlepage\end{frame}

\begin{frame}{``My program is slow'': but how slow, and for what input?}
\begin{columns}[T]
\begin{column}{0.5\textwidth}
{\small
You time two programs: A takes 2 seconds, B takes 5 seconds. Is A better?

\vspace{2mm}
Not necessarily:
\begin{itemize}\setlength\itemsep{2pt}
\item the timings depend on \textbf{your computer}, other programs, the Python version;
\item they say nothing about a \textbf{bigger input}. A might take 2 s for 1\,000 items but 3 hours for 100\,000.
\end{itemize}
\vspace{1mm}
So computer scientists ask a different question: \key{how does the number of steps grow as the input grows?}}
\end{column}
\begin{column}{0.46\textwidth}
\begin{block}{After this lecture you can}
{\small
\begin{itemize}\setlength\itemsep{2pt}
\item count the steps of a loop or a function;
\item recognise $O(1)$, $O(\log n)$, $O(n)$, $O(n\log n)$, $O(n^2)$ and $O(2^n)$;
\item simplify an expression to its Big-O form;
\item read the complexity off a piece of code;
\item see why the difference matters, with real timings.
\end{itemize}}
\end{block}
\end{column}
\end{columns}
\end{frame}

\begin{frame}{Constant and linear: counting steps}
\begin{columns}[T]
\begin{column}{0.52\textwidth}
\pycode{constant}
\vspace{2mm}
\pycode{linear}
\end{column}
\begin{column}{0.44\textwidth}
{\small
\textbf{\texttt{first\_item}}: one step, no matter how long the list is. We call this \key{constant time}, $O(1)$.

\vspace{3mm}
\textbf{\texttt{total}}: the loop body runs once per item, so the work is \emph{proportional} to $n$. Double the list, double the time: \key{linear time}, $O(n)$.

\vspace{3mm}
We count the \textbf{basic operations} (an addition, a comparison, an assignment), not seconds. Whether a step is $1$ or $3$ machine instructions does not change how the \emph{growth} looks.}
\end{column}
\end{columns}
\end{frame}

\begin{frame}{Nested loops: quadratic time}
\begin{columns}[T]
\begin{column}{0.52\textwidth}
\pycode{quadratic}
\vspace{2mm}
\pyout{pairs}
\end{column}
\begin{column}{0.44\textwidth}
{\small
The inner loop runs about $n$ times for each of the $n$ values of \texttt{i}. With no duplicate the function compares every pair once:
\[
\binom{n}{2}=\frac{n(n-1)}{2}=\frac{n^2}{2}-\frac{n}{2}.
\]
For large $n$ the $n^2$ term dominates, so this is \key{quadratic time}, $O(n^2)$.

\vspace{2mm}
\textbf{Rule of thumb:} two nested loops over the input $\Rightarrow$ $O(n^2)$; three $\Rightarrow$ $O(n^3)$.}
\end{column}
\end{columns}
\end{frame}

\begin{frame}{Logarithmic time: halving}
\begin{columns}[T]
\begin{column}{0.46\textwidth}
\pycode{logarithmic}
\vspace{2mm}
{\small
\texttt{halvings(1000)} $=9$, \texttt{halvings(1000000)} $=19$. How many times can you halve $n$ before reaching 1? About $\log_2 n$ times.

\vspace{2mm}
\key{Logarithmic time}, $O(\log n)$: a million items need ${\sim}20$ steps, a billion only ${\sim}30$. Binary search is the classic example; so is walking down a balanced tree.}
\end{column}
\begin{column}{0.5\textwidth}
\pyout{steps}
\vspace{2mm}
{\small Compare the columns for $n=100\,000$: $\log_2 n$ is \textbf{16}, $n$ is 100\,000, $n\log_2 n$ is 1.6 million, and $n^2$ is \textbf{ten billion}.}
\end{column}
\end{columns}
\end{frame}

\begin{frame}{The common growth rates}
\begin{columns}[T]
\begin{column}{0.44\textwidth}
\renewcommand{\arraystretch}{1.28}\small
\begin{tabular}{@{}lll@{}}
\toprule
Big-O & Name & Example\\\midrule
$O(1)$ & constant & \texttt{items[i]}\\
$O(\log n)$ & logarithmic & binary search\\
$O(n)$ & linear & scan a list\\
$O(n\log n)$ & linearithmic & merge sort\\
$O(n^2)$ & quadratic & nested loops\\
$O(2^n)$ & exponential & all subsets\\
\bottomrule
\end{tabular}
\vspace{2mm}

{\small Slower-growing is better: $O(1)\prec O(\log n)\prec O(n)\prec O(n\log n)\prec O(n^2)\prec O(2^n)$.}
\end{column}
\begin{column}{0.54\textwidth}
\begin{sfig}[0.78]
\begin{tikzpicture}
\begin{axis}[width=9.4cm,height=6.6cm,xmin=0,xmax=20,ymin=0,ymax=100,xlabel={input size $n$},ylabel={steps},axis lines=left,legend pos=north west,legend style={font=\scriptsize},tick label style={font=\scriptsize},label style={font=\small},clip=true]
  \addplot[tdGreen,very thick,domain=1:20] {1.5};
  \addplot[tdBlue!70!tdGreen,very thick,domain=1:20,samples=60] {4*ln(x)/ln(2)/1.5};
  \addplot[tdBlue,very thick,domain=1:20] {x*2.2};
  \addplot[tdPurple,very thick,domain=1:20,samples=60] {x*ln(x)/ln(2)/1.55};
  \addplot[tdAmber,very thick,domain=1:20,samples=60] {x*x/4};
  \addplot[tdRed,very thick,domain=1:20,samples=60] {2^x/9000*3};
  \legend{$1$,$\log n$,$n$,$n\log n$,$n^2$,$2^n$}
\end{axis}
\end{tikzpicture}
\end{sfig}
{\tiny curves rescaled to fit one picture: only the \emph{shape} matters.}
\end{column}
\end{columns}
\end{frame}

\begin{frame}{Big-O: dropping what does not matter}
\begin{columns}[T]
\begin{column}{0.5\textwidth}
\begin{block}{Informal definition}
{\small $f(n)=O(g(n))$ means: for large enough $n$, $f$ grows \textbf{no faster than a constant times $g$}. Big-O is an \emph{upper bound} on growth.}
\end{block}
\vspace{2mm}
{\small
\textbf{Two simplification rules}
\begin{enumerate}\setlength\itemsep{2pt}
\item \textbf{Drop constant factors}: $5n$ and $100n$ are both $O(n)$.
\item \textbf{Keep only the fastest-growing term}: $n^2+n$ is $O(n^2)$.
\end{enumerate}}
\end{column}
\begin{column}{0.46\textwidth}
\renewcommand{\arraystretch}{1.3}\small
\begin{tabular}{@{}ll@{}}
\toprule
Steps counted & Big-O\\\midrule
$7$ & $O(1)$\\
$3n+5$ & $O(n)$\\
$\tfrac12 n^2-\tfrac12 n$ & $O(n^2)$\\
$n^2+1000n+10^6$ & $O(n^2)$\\
$5n\log_2 n+20n$ & $O(n\log n)$\\
$2^n+n^{10}$ & $O(2^n)$\\
\bottomrule
\end{tabular}
\vspace{3mm}

{\small \textbf{Why drop terms?} For $n=1000$, in $n^2+1000n$ the $n^2$ part is already 1\,000\,000 and the other part only 1\,000.}
\end{column}
\end{columns}
\end{frame}

\begin{frame}{Quiz: what is the Big-O of each function?}
\begin{columns}[T]
\begin{column}{0.48\textwidth}
\pycode{quiz_a}
\vspace{2mm}
\pycode{quiz_b}
\end{column}
\begin{column}{0.48\textwidth}
\pycode{quiz_c}
\vspace{2mm}
\pycode{quiz_d}
\end{column}
\end{columns}
\vspace{2mm}
{\small Think first, then turn the page. Hint: look at how many times each loop runs, and whether a loop variable is added to, or multiplied or divided.}
\end{frame}

\begin{frame}{Quiz: answers}
\begin{columns}[T]
\begin{column}{0.5\textwidth}
\pyout{quiz}
\vspace{2mm}
{\small Doubling $n$ from 1000 to 2000 shows the growth: $\times2$, $+1$, $\times4$, and a bit more than $\times2$.}
\end{column}
\begin{column}{0.46\textwidth}
{\small
\begin{description}\setlength\itemsep{3pt}
\item[\texttt{mystery\_a}] $5n$ steps: the inner loop is a \emph{fixed} 5 times. $\Rightarrow O(n)$.
\item[\texttt{mystery\_b}] \texttt{i} doubles each time, so about $\log_2 n$ rounds. $\Rightarrow O(\log n)$.
\item[\texttt{mystery\_c}] two nested loops of $n$: $n^2$ steps. $\Rightarrow O(n^2)$.
\item[\texttt{mystery\_d}] $n$ rounds of an inner loop that halves: $n\log_2 n$. $\Rightarrow O(n\log n)$.
\end{description}}
\end{column}
\end{columns}
\end{frame}

\begin{frame}{Exponential time: when the recursion repeats itself}
\begin{columns}[T]
\begin{column}{0.46\textwidth}
\pycode{fib_calls}
\vspace{2mm}
\pyout{fib_demo}
\vspace{1mm}
{\small Each call makes \emph{two} more calls, so the work roughly multiplies by $1.6$ for every $+1$ in $n$. That is $O(2^n)$-like growth.}
\end{column}
\begin{column}{0.5\textwidth}
\begin{sfig}[0.45]
\begin{tikzpicture}
\begin{axis}[width=9cm,height=5.4cm,ymode=log,xlabel={$n$},ylabel={function calls},axis lines=left,tick label style={font=\scriptsize},label style={font=\small}]
  \addplot[tdRed,very thick,mark=*,mark size=1.6pt] table[x=n,y=calls]{snippets/fib_data.out};
\end{axis}
\end{tikzpicture}
\end{sfig}
{\tiny A straight line on a log axis means exponential growth.}
\vspace{2mm}

\pyout{growth}
\end{column}
\end{columns}
\end{frame}

\begin{frame}{The difference you can feel}
\begin{columns}[T]
\begin{column}{0.5\textwidth}
{\footnotesize\textbf{Is $x$ in this collection?} for 100\,000 items:}
\pyout{membership}
\vspace{2mm}
{\small A list scans item by item, $O(n)$. A \texttt{set} hashes the value and jumps straight to it: about $O(1)$.}
\end{column}
\begin{column}{0.46\textwidth}
{\footnotesize\textbf{Building a list of $m$ numbers}, two ways:}
\pyout{append_insert}
\vspace{2mm}
{\small \texttt{append} adds at the end: $O(1)$ each, so $O(m)$ in total. \texttt{insert(0, x)} must shift every existing item: $O(m)$ each, $O(m^2)$ in total. Doubling $m$ makes it about four times slower.}
\end{column}
\end{columns}
\end{frame}

\begin{frame}{Best, worst, average, and memory}
\begin{columns}[T]
\begin{column}{0.5\textwidth}
\begin{block}{Which case do we report?}
{\small
Take linear search:
\begin{itemize}\setlength\itemsep{2pt}
\item \textbf{best} case: target is first, $O(1)$;
\item \textbf{worst} case: last or absent, $O(n)$;
\item \textbf{average}: about $n/2$ steps, still $O(n)$.
\end{itemize}
Unless told otherwise, quote the \textbf{worst case}: it is a promise that holds for \emph{every} input.}
\end{block}
\end{column}
\begin{column}{0.46\textwidth}
\begin{block}{Time is not the only resource}
{\small
\textbf{Space complexity} counts extra memory.\\
\texttt{total(items)} uses $O(1)$ extra space; \texttt{merge\_sort} builds new lists, $O(n)$ extra space.

\vspace{2mm}
Often there is a \textbf{trade-off}: a \texttt{set} answers lookups fast, but spends memory to do it.

\vspace{2mm}
\textbf{Amortised} cost: \texttt{list.append} is usually fast and occasionally copies the list when it grows, but averaged over many appends it is $O(1)$.}
\end{block}
\end{column}
\end{columns}
\end{frame}

\begin{frame}{Cost of common Python operations}
\renewcommand{\arraystretch}{1.25}\small
\begin{tabularx}{\textwidth}{@{}>{\raggedright}p{54mm}>{\raggedright}p{30mm}X@{}}
\toprule
Operation & Cost & Note\\\midrule
\texttt{len(items)},\ \texttt{items[i]} & $O(1)$ & direct access by position\\
\texttt{items.append(x)} & $O(1)$ amortised & adds at the end\\
\texttt{items.insert(0, x)},\ \texttt{items.pop(0)} & $O(n)$ & shifts every item\\
\texttt{x in items} (list) & $O(n)$ & linear search\\
\texttt{x in s}, \texttt{s.add(x)} (set, dict) & $O(1)$ on average & hashing\\
\texttt{sorted(items)} & $O(n\log n)$ & Timsort\\
\texttt{for x in items:} & $O(n)$ & one pass\\
\bottomrule
\end{tabularx}
\vspace{3mm}

{\small A nested loop \emph{hides} cost: \texttt{for x in a: if x in b:} is $O(n\cdot m)$ when \texttt{b} is a list, but only $O(n)$ when \texttt{b} is a set.}
\end{frame}

\begin{frame}{Recap and practice}
\begin{columns}[T]
\begin{column}{0.48\textwidth}
\begin{block}{Take-aways}
{\small
\begin{enumerate}\setlength\itemsep{3pt}
\item Big-O describes \key{how steps grow with $n$}, not seconds.
\item Drop constants and lower-order terms.
\item Common classes, slow to fast growth: $2^n$, $n^2$, $n\log n$, $n$, $\log n$, $1$.
\item One loop $O(n)$; nested loops $O(n^2)$; halving $O(\log n)$; halving inside a loop $O(n\log n)$.
\item The choice of data structure (list vs.\ set) can change the Big-O.
\end{enumerate}}
\end{block}
\end{column}
\begin{column}{0.48\textwidth}
\begin{block}{Practice}
{\small
\begin{enumerate}\setlength\itemsep{3pt}
\item State the Big-O of \texttt{sum(range(n))}, \texttt{max(items)} and \texttt{items.count(x)}.
\item Write a function that finds a pair of numbers in a list that adds to a target: first in $O(n^2)$, then in $O(n)$ using a \texttt{set}.
\item \texttt{for i in range(n): for j in range(i):} How many times does the body run? Is it still $O(n^2)$?
\item Use \texttt{timeit} to check that \texttt{x in list} doubles when the list doubles.
\item Why is $O(n^2)$ acceptable for $n=100$ but dangerous for $n=10^6$?
\end{enumerate}}
\end{block}
\end{column}
\end{columns}
\end{frame}
\end{document}
